\documentclass[12pt]{article}

% Idioma e codificação
\usepackage[brazil]{babel}
\usepackage[utf8]{inputenc}
\usepackage[T1]{fontenc}

% Matemática e listas
\usepackage{amsmath, amssymb}
\usepackage{enumitem}
\usepackage{float}
\usepackage{colortbl}
\usepackage{geometry}
\usepackage{graphicx}
\usepackage{xcolor}

% Cores institucionais e personalizadas
\definecolor{maincolor}{RGB}{0,100,50}
\definecolor{forestgreen}{RGB}{0,100,50}
\definecolor{bgsql}{RGB}{245, 245, 245}
\definecolor{verdeuni}{RGB}{0, 100, 0}
\definecolor{blueuni}{RGB}{0, 0, 128}

% Cabeçalho e rodapé
\usepackage{fancyhdr}
\usepackage{tabularx}
\usepackage{booktabs}
\usepackage[most]{tcolorbox}
\tcbuselibrary{minted}

\newtcblisting{sqlcode}{
    listing engine=minted,
    minted language=sql,
    colback=bgsql,
    colframe=verdeuni,
    listing only,
    title=SQL,
    fonttitle=\bfseries\color{white},
    coltitle=white, 
    colbacktitle=verdeuni
}

\geometry{margin=2.5cm}
\pagestyle{fancy}
\fancyhf{}
\renewcommand{\headrulewidth}{2pt}
\renewcommand{\footrulewidth}{2pt}
\renewcommand{\headrule}{\hbox to\headwidth{\color{forestgreen}\leaders\hrule height \headrulewidth\hfill}}
\renewcommand{\footrule}{\hbox to\headwidth{\color{forestgreen}\leaders\hrule height \footrulewidth\hfill}}
\setlength{\headheight}{52pt}

\fancyhead[L]{\textbf{UniRV}}
\fancyhead[R]{%
  \raggedleft
  \textcolor{forestgreen}{\small \textbf{Lista de Exercícios Complementar}\\Me.~Sergio Souza Novak\\\textit{Banco de Dados}}
}
\fancyfoot[C]{\textcolor{forestgreen}{\thepage}}

\begin{document}

\section*{Lista de Exercícios – Criação, Manipulação e Análise de Dados}

\begin{tcolorbox}[colback=blue!5,colframe=blue!40!black,title=Objetivo da Atividade]
Nesta atividade, você deverá realizar o ciclo completo: criar a estrutura das tabelas (DDL), inserir dados de teste (DML) e realizar consultas analíticas (DQL) utilizando filtros e agrupamentos.
\end{tcolorbox}

\hrulefill

\section*{Exercício 1 – Gestão de Estoque (DDL, Insert e Select)}



\textbf{Cenário:} Modele uma tabela de \texttt{produto} com: ID (serial), nome (obrigatório), descrição, preço (decimal), estoque (default 0), SKU (único, 8 caracteres) e data automática.

\textbf{Tarefas:}
\begin{enumerate}
    \item Escreva o comando \texttt{CREATE TABLE}.
    \item Insira 3 produtos: 'Teclado Mecânico' (250.00, 15 un), 'Mouse Gamer' (150.00, 5 un) e 'Monitor 24pol' (890.00, 0 un).
    \item Crie uma consulta que retorne nome e preço dos produtos com \textbf{estoque > 0}, ordenados do mais caro ao mais barato.
\end{enumerate}

\hrulefill

\section*{Exercício 2 – Relatório de Vendas (Group By e Count)}



\textbf{Cenário:} Considere a tabela \texttt{venda} com: \texttt{id}, \texttt{data\_venda}, \texttt{valor} e \texttt{vendedor\_id}.

\textbf{Tarefas:}
\begin{enumerate}
    \item Escreva a query que mostre a quantidade de vendas realizada por cada vendedor.
    \item Utilize \texttt{HAVING} para mostrar apenas vendedores com \textbf{mais de 10 vendas}.
\end{enumerate}

\hrulefill

\section*{Exercício 3 – Status de Clientes (DDL e Having)}



\textbf{Cenário:} Tabela \texttt{cliente} com ID, nome, CPF (único), e-mail e um campo \texttt{ativo} (booleano).

\textbf{Tarefas:}
\begin{enumerate}
    \item Crie a tabela e insira 2 clientes ativos e 1 inativo.
    \item Conte quantos clientes existem para cada status (\texttt{ativo}).
    \item Exiba apenas o status que possui \textbf{menos de 2 clientes} cadastrados.
\end{enumerate}

\hrulefill

\section*{Exercício 4 – Filtros Compostos em Frota}



\textbf{Cenário:} Tabela \texttt{carro} com placa (única), marca, modelo, ano e valor.

\textbf{Tarefas:}
\begin{enumerate}
    \item Escreva um \texttt{SELECT} para contar quantos carros da marca 'Toyota' foram fabricados \textbf{após 2020}.
    \item Agrupe a contagem de carros por marca e ordene o resultado alfabeticamente.
\end{enumerate}

\hrulefill

\section*{Exercício 5 – Folha de Pagamento (Sum e Where)}



\textbf{Cenário:} Tabela \texttt{empregado} com nome, cargo, salário e status ativo.

\textbf{Tarefas:}
\begin{enumerate}
    \item Calcule a \textbf{soma total} dos salários apenas de funcionários \textbf{ativos}.
    \item Explique por que usamos \texttt{SUM(salario)} e não \texttt{COUNT(salario)} para obter o custo total da folha.
\end{enumerate}

\hrulefill

\section*{Desafio de Lógica – Correção de Sintaxe}

Considere o código abaixo que apresenta erros de ordem de execução:

\begin{sqlcode}
SELECT departamento, COUNT(id)
FROM funcionarios
HAVING COUNT(id) > 5
WHERE salario > 2000
GROUP BY departamento;
\end{sqlcode}

\textbf{Tarefa:}
\begin{enumerate}
    \item Corrija o código respeitando a ordem: \texttt{FROM} $\rightarrow$ \texttt{WHERE} $\rightarrow$ \texttt{GROUP BY} $\rightarrow$ \texttt{HAVING}.
    \item Explique por que o \texttt{WHERE} deve vir antes do \texttt{GROUP BY}.
\end{enumerate}



\begin{tcolorbox}[colback=green!5,colframe=green!40!black,title=Dica Final]
Sempre visualize o diagrama de classes antes de codificar. Se o atributo é único no diagrama, ele deve ser \texttt{UNIQUE} no SQL!
\end{tcolorbox}

\end{document}